 \section{Background}\label{sec:background}
\subsection{Defining Basic Terms}\label{subsec:defining-basic-terms}

    Before beginning the analysis, we will define some terms and concepts used throughout the remainder of the paper.


    \subsubsection{First-Order Separable ODEs}

    A separable ODE has the form
    \[
        \frac{dy}{dx}=f(x)g(y).
    \]
    It is considered separable as it can be rearranged to \textit{separate} x and y onto opposite sides of the equation:
    \[
        \int\frac{1}{g(y)}dy=\int f(x)dx.
    \]
    When both sides are integrated, the resulting solution relates the two variables, plus the constant of integration. The original differential equation describes a family of curves using their slopes, while a general solution describes a family of curves related by the constant of integration, and a specific solution represents the particular curve which satisfies a set of initial conditions \parencite{MathematicsAnalysis2021}.

    \subsubsection{Second-Order Linear ODEs with Constant Coefficients}

    Given in the form
    \[
        ay''+by'+cy=0,
    \]
    We can make an educated guess about what $y$ may be. Since differentiating an exponential is equivalent to multiplying by a constant, a probable solution is $y=e^{rx}$. Then $y'=re^{rx}$, $y''=r^2 e^{rx}$. Substituting this in gives
    \[
        ar^2 e^{rx}+bre^{rx}+ce^{rx}=0 \Longrightarrow e^{rx}\left(ar^2+br+c \right)=0.
    \]
    As $e^{rx}$ can never equal zero, this forces the characteristic equation:
    \[
        ar^2+br+c=0.
    \]
    When we solve this equation for $r$, three potential types of solutions for $e^{rx}$ exist.
    \begin{enumerate}
        \item In the first case, there will be two real and distinct roots $r_1$, $r_2$.
        This gives the general solution
        $y=Ae^{r_1 x}+Be^{r_2 x}$.
        \item In the second case, the real root $r$ is repeated.
        This gives the general solution $y=(A+Bx)e^{rx}$.

        \item In the final case, the roots are complex conjugates in the form $r=\alpha\pm i\beta$, which gives the solution $y=C_1 e^{(\alpha+ i\beta)x}+C_2 e^{(\alpha- i\beta)x}$.
        Using Euler's formula, we can convert this into a real-valued function:
        \[
            y=e^{\alpha x}(A\cos\beta x+B\sin\beta x).
        \]
        In terms of the physical system, the appearance of complex roots causes changes to the oscillation of the system.
        The behaviours are listed in the table below:

    \end{enumerate}

    \begin{table}[H]
        \centering
        \begin{tabular}{ll}
            \toprule
            \textbf{Value of $\alpha$} & \textbf{Behavior} \\
            \midrule
            $\alpha > 0$ & growing oscillation (unstable) \\
            $\alpha = 0$ & perfect harmonic motion (stable) \\
            $\alpha < 0$ & decaying oscillation (stable) \\
            \bottomrule
        \end{tabular}
        \caption{Value of $\alpha$ and Corresponding Behavior}
        \label{tab:alphabehavior}
    \end{table}

    \subsubsection{Euler's Method}

    Euler's method provides a numerical approximation to a solution of $y'=f(x,y)$ when it cannot be solved analytically. 
    The method accomplishes this by stepping along the tangent line at a given point:
    \[
        y_{n+1}=y_n+hf(x_n,y_n), \qquad x_{n+1} = x_n+h.
    \]
    Where $h$ is the step size \parencite{MathematicsAnalysis2021}.
    Obviously, the error increases as $h$ increases.

    \subsubsection{Maclaurin Series}

    The Maclaurin series is simply a Taylor series centred at the origin ($x=0$). It expresses a function in terms of an infinite polynomial created from its derivatives:
    \[
        f(x)=f(0)+f'(0)x+\frac{f''(0)}{2!}x^2+\frac{f'''(0)}{3!}x^3+\dots=\sum_{n=0}^\infty\frac{f^{(n)}(0)}{n!}x^n
    \]
    \parencite{MathematicsAnalysis2021}.
    This formula is necessary for linearizing around an equilibrium. Using this
    series, we can approximate solutions more accurately near $x=0$.